\section{Introduction}

Modern AI systems can combine models, agents, or other learners with different quality, cost, latency, competence, and learning profiles. Their immediate operational question is familiar: which member should handle the current task? If members can learn, however, this view is incomplete. Operation can generate experience; capable interventions can generate potentially reusable knowledge; and learning can change future competence.

The programme-level motivation is a feedback loop:
\begin{equation}
\begin{gathered}
\text{heterogeneous learners} \rightarrow \text{operation and routing} \rightarrow \text{experience}\\
\rightarrow \text{learning or transfer} \rightarrow \text{competence redistribution}\\
\rightarrow \text{future operation}.
\end{gathered}
\label{eq:cycle}
\end{equation}
The system consequently faces two linked allocation problems: work allocation and knowledge or learning allocation. Together they determine a future organization of competence, rather than only use a current one. We use a provisional hierarchical and collaborative language of inexpensive learners and capable teachers: teachers can provide service and potentially useful knowledge, while learners can absorb demand and adapt. These are functional roles, not fixed types.

Environmental dynamics make the distinction especially relevant. A persistent change in demand, costs, latency requirements, resources, constraints, or required competences can make a previously useful organization less suitable. Whether an HLS can deliberately reorganize its competences to create long-term value is a hypothesis, not an established response to change.

The repository-level research question is:
\begin{quote}
Can the deliberate evolution of competence allocation in a heterogeneous learning system improve its long-term performance compared with independently optimizing task allocation and knowledge transfer?
\end{quote}
This draft does not replace RQ0.

Closed-loop routing and learning already have direct operational precedent. For example, RouteNLP uses routing failures to select targeted distillation data and subsequently recalibrates routing \citep{guo2026routenlp}. Cross-domain models also integrate work assignment, competence dynamics, and training \citep{borgonjon2024dynamic}. These precedents constrain component claims; they do not by themselves settle the programme-level question.

The purpose of the paper is conservative. It documents a provisional HLS architecture, an integration-theory scaffold, and preliminary minimal analytical results with numerical consistency checks. The architecture and general theory are still being developed; general experiments do not yet exist. The results therefore do not establish novelty, general HLS superiority, realistic empirical validation, or the candidate coupling proposition P4.
